\documentclass[10pt,a4paper]{article}

% ==========================================
% Minimalist Academic Style & Packages
% ==========================================
\usepackage[utf8]{utf8}
\usepackage[margin=1.0in]{geometry}
\usepackage{amsmath,amssymb,amsfonts}
\usepackage{microtype}
\usepackage{hyperref}
\usepackage{xcolor}
\usepackage{graphicx}
\usepackage{booktabs}
\usepackage{caption}
\usepackage{enumitem}
\usepackage{titlesec}
\usepackage{listings}

% Brand Accent Colors (Google Research Minimalist Style)
\definecolor{googleblue}{HTML}{1A73E8}
\definecolor{googletext}{HTML}{202124}
\definecolor{googlecodebg}{HTML}{F8F9FA}
\definecolor{googlecomment}{HTML}{5F6368}

% Hyperlink Formatting
\hypersetup{
  colorlinks=true,
  linkcolor=googleblue,
  citecolor=googleblue,
  urlcolor=googleblue
}

% Section Titles Styling
\titleformat{\section}
  {\color{googletext}\large\bfseries}
  {\thesection}{1em}{}
  [\color{googleblue}\titlerule]

\titleformat{\subsection}
  {\color{googletext}\normalfont\bfseries}
  {\thesubsection}{1em}{}

% Code Block Styling
\lstset{
  backgroundcolor=\color{googlecodebg},
  basicstyle=\ttfamily\small,
  breaklines=true,
  captionpos=b,
  commentstyle=\color{googlecomment},
  keywordstyle=\color{googleblue}\bfseries,
  stringstyle=\color{googletext},
  frame=single,
  rulecolor=\color{googlecodebg}
}

\captionsetup{labelfont={bf,color=googletext},textfont={color=googletext},skip=6pt}

\color{googletext}

% ==========================================
% Metadata
% ==========================================
\title{
  \vspace{-1.5cm}
  {\color{googleblue}\rule{\textwidth}{2pt}}\\[0.4cm]
  \Large\bfseries Appendix A: Porting the Autonomous Cognitive Control Plane to Google Vertex AI\\
  \large Cross-Cloud Invariants, XLA Compilation Semantics, and Substrate Translation
  \vspace{0.2cm}
}

\author{
  \textbf{Paulo H. Leocadio}\\
  \small\textit{Author, "Thinking Cybersecurity: Cognitive Control Planes and Autonomous Defense Systems"}\\
  \small\texttt{research@domain.org}
}

\date{\small\today}

% ==========================================
% Body
% ==========================================
\begin{document}

\maketitle

\begin{abstract}
\noindent
Autonomous cybersecurity systems do not fail because models are inaccurate; they fail because deployment architectures are fragile. In this appendix, we extend the canonical control-plane architecture established in Chapter 10---originally operationalized on AWS SageMaker and PyTorch---to Google Cloud Platform (GCP). We stress-test the five foundational architectural invariants under cross-cloud translation: container discipline, runtime substrate stability, registry/rollback semantics, identity containment, and agentic orchestration compatibility. Furthermore, we analyze the computational divergence introduced when transitioning tensor execution from direct CUDA kernels to PyTorch/XLA compilation on Tensor Processing Units (TPUs), demonstrating how long-context ingestion on Vertex AI compresses system response latency ($R_s$) without compromising policy gating guarantees.
\end{abstract}

\vspace{0.3cm}

\section{Introduction \& Research Context}
Throughout \textit{Thinking Cybersecurity}, Hugging Face Diffusers and open transformer-based models serve as modular, composable primitives. However, infrastructure appears only when these primitives integrate into training pipelines, inference endpoints, container orchestration, and policy gating loops.

While Chapter 10 established the canonical stack using PyTorch and Amazon SageMaker as the reference orchestration substrate, the underlying control plane is cloud-agnostic. In this appendix, we execute a formal architectural translation:
\begin{equation}
  \mathcal{T}(\mathcal{C}_{\text{AWS}}) \rightarrow \mathcal{C}_{\text{GCP}}
\end{equation}
where $\mathcal{C}_{\text{AWS}}$ represents the reference AWS implementation and $\mathcal{C}_{\text{GCP}}$ represents the target Google Cloud environment, ensuring that all structural invariants remain preserved.

\section{Architectural Invariants & Cross-Cloud Translation}
Portability is evaluated not by superficial API feature equivalence, but by structural invariance. Across all target environments, five non-negotiable invariants must hold:
\begin{enumerate}[leftmargin=1.5em, itemsep=0.2em]
  \item \textbf{Container Discipline:} Deterministic, container-native execution of custom PyTorch-based pipelines with absolute environment reproducibility.
  \item \textbf{Runtime Substrate Stability:} Numerical predictability within bounded tolerance ($\epsilon$) under adversarial load.
  \item \textbf{Registry and Rollback Semantics:} Traceable model lineage and deterministic rollback without infrastructure drift.
  \item \textbf{Governance and Identity Containment:} Least-privilege IAM boundaries across training, inference, and tool invocation layers.
  \item \textbf{Agentic Orchestration Compatibility:} Deterministic execution of the iterative control loop:
  \begin{equation}
    S_{t+1} = \mathcal{F}(S_t, O_t, \pi)
  \end{equation}
  where $S_t$ is system state, $O_t$ is observation, and $\pi$ represents policy constraints.
\end{enumerate}

\section{Execution Substrate Divergence: CUDA vs. TPU / XLA}
The primary runtime divergence on GCP occurs at the hardware compilation layer. While AWS executes PyTorch workloads directly via NVIDIA CUDA kernels, Google Cloud introduces Tensor Processing Units (TPUs) via Accelerated Linear Algebra (XLA) compilation:
\begin{align*}
  \text{AWS Pathway:} \quad &\text{PyTorch (CUDA)} \rightarrow \text{Direct Tensor Execution} \\
  \text{GCP Pathway:} \quad &\text{PyTorch (TPU)} \rightarrow \text{XLA Compilation} \rightarrow \text{Tensor Execution}
\end{align*}

For iterative denoising and attention calculations in security diffusion pipelines ($x_{t-1} = \alpha_t x_t - \beta_t \epsilon_\theta(x_t, t)$), XLA graph optimizations (operation fusion and reordering) must be verified to ensure numerical approximation error remains within operational tolerance:
\begin{equation}
  \| f(x; W) - f(x; W_q) \| \le \epsilon
\end{equation}

\section{Code Scaffold: Vertex AI Deployment & Policy Gate}
The following listing demonstrates the translation of the canonical deployment loop to Vertex AI, incorporating the mandatory Policy Gate ($\pi$) middleware:

\subsection{Embedded Policy Container Application}
This web application runs inside the custom managed prediction container (e.g., \texttt{app.py}). It exposes the required HTTP endpoints to conform with Vertex AI's container specifications while dynamically validating model inferences against a designated security policy.

\begin{lstlisting}[language=Python, caption={Custom Embedded Policy Gate Container Application (app.py).}]
from fastapi import FastAPI, Request
import json

app = FastAPI()

class SecurityPolicyGate:
    """Enforces policy constraints for state transitions and safety."""
    def __init__(self, restricted_actions):
        self.forbidden = restricted_actions
        
    def validate(self, inference_result):
        # 1. Content Safety & Input Sanitization
        if any(trigger in inference_result.lower() for trigger in self.forbidden):
            return False
        # 2. Structural Constraint: Event Horizon Check
        if "action" not in inference_result.lower():
            return False
        return True

# Initialize the runtime gate with forbidden triggers
gate = SecurityPolicyGate(restricted_actions=["malware", "exploit"])

@app.get("/healthz")
async def healthz():
    """Required by Vertex AI to continuously verify container health."""
    return {"status": "healthy"}

@app.post("/predict")
async def predict(request: Request):
    """Intercepts and sanitizes predictions generated by the underlying model."""
    body = await request.json()
    
    # Run prediction logic using your underlying model artifact
    # (Placeholder logic assuming a standard string return)
    inference_result = "action: executed benign sequence" 
    
    # Enforce security constraints via the policy gate
    if not gate.validate(inference_result):
        return {"predictions": [{"error": "Policy violation detected", "status": "REJECTED"}]}
        
    return {"predictions": [inference_result]}
\end{lstlisting}

\subsection{Vertex AI Orchestration Pipeline}
This script (e.g., \texttt{deploy.py}) targets the Google Cloud Platform environment to register the model image and provision the managed cloud infrastructure.

\begin{lstlisting}[language=Python, caption={Vertex AI Endpoint Deployment Script (deploy.py).}]
import google.cloud.aiplatform as aiplatform

def deploy_canonical_to_vertex(project_id, location, container_uri):
    """Uploads the custom policy container and provisions a managed endpoint."""
    # Initialize the Vertex AI SDK context
    aiplatform.init(project=project_id, location=location)
    
    # 1. Register Model in Vertex AI Model Registry pointing to your container
    # Explicitly map the container paths to the FastAPI application routes
    model = aiplatform.Model.upload(
        display_name="cyber-cognitive-cannon-v1",
        serving_container_image_uri=container_uri,
        serving_container_predict_route="/predict",
        serving_container_health_route="/healthz"
    )
    
    # 2. Deploy to a Managed Hardware Endpoint
    endpoint = model.deploy(
        machine_type="n1-standard-4",
        accelerator_type="NVIDIA_TESLA_T4",
        accelerator_count=1,
        min_replica_count=1,
        max_replica_count=2
    )
    return endpoint
\end{lstlisting}


\section{Conclusion}
The comparative analysis confirms that the canonical cognitive architecture is portable across cloud environments when architectural invariants are maintained. Google Cloud's Vertex AI and long-context capabilities compress latency and enhance state estimation, while GKE and Cloud IAM provide the necessary containment boundaries. The canon remains architectural; the cloud substrate remains contextual.

\begin{thebibliography}{9}
\bibitem{leocadio2026}
P.~H. Leocadio, \emph{Thinking Cybersecurity: Cognitive Control Planes and Autonomous Defense Systems}, Chapter 10, 2026.
\end{thebibliography}

\end{document}